\documentclass[11pt,a4paper]{article}

\usepackage[utf8]{inputenc}
\usepackage[T1]{fontenc}
\usepackage{lmodern}
\usepackage{amsmath,amssymb,amsthm}
\usepackage{booktabs}
\usepackage{enumitem}
\usepackage{hyperref}
\usepackage[margin=2.5cm]{geometry}
\usepackage{xcolor}

\definecolor{codeblue}{RGB}{30,90,156}

\newtheorem{theorem}{Theorem}[section]
\newtheorem{proposition}[theorem]{Proposition}
\newtheorem{corollary}[theorem]{Corollary}
\newtheorem{lemma}[theorem]{Lemma}

\theoremstyle{definition}
\newtheorem{definition}[theorem]{Definition}

\theoremstyle{remark}
\newtheorem*{remark}{Remark}

\hypersetup{
  colorlinks=true,
  linkcolor=codeblue,
  citecolor=codeblue,
  urlcolor=codeblue,
}

\newcommand{\Jcost}{J}

\title{%
  One Involution:\\
  The Recognition Cost and the Golden Ratio\\
  from a Single Symmetry}
\author{%
  Jonathan Washburn\\
  Recognition Physics Institute, Austin, TX, USA\\
  \texttt{jon@recognitionphysics.org}}
\date{May 2026}

\begin{document}

\maketitle

\begin{abstract}\sloppy
The Recognition Science program treats two quantities as fundamental: the
recognition cost $J(x)=\tfrac12(x+x^{-1})-1$, and the golden ratio $\varphi$ that
sets the scale of its self-similar hierarchy. These look like two independent
inputs. They are not. Both are determined by a single object, the reciprocal
involution $\iota(x)=x^{-1}$ on the positive reals. The cost is exactly the
$\iota$-symmetric cost, and the fixed point of $\iota$ is precisely the cost's
unique zero, the unit. The golden ratio is exactly the unique fixed point greater
than one of the affine shift $1+\iota$, the map obtained by perturbing $\iota$ by
its own fixed point. So the unit and the scale, the two constants the framework
is built from, are the two fixed points of one involution family $\{\iota,\,
1+\iota\}$. We state this precisely, prove it, and are explicit about what it does
and does not claim. Every assertion is machine-checked in Lean~4 against Mathlib,
with no \texttt{sorry} and no Recognition-Science-specific axiom; the module is
\texttt{IndisputableMonolith.Foundation.UniversalForcing.ReciprocalGenerator}.
\end{abstract}

\section{The claim, and why it is not obvious}

A reader who has seen the cost $J$ and the golden ratio $\varphi$ derived
separately in the Recognition Science literature could reasonably take them to be
two distinct primitives: one a function pinned by a functional equation, the other
a number pinned by a self-similarity constraint. The cost-uniqueness theorem fixes
$J$ from reciprocal symmetry, normalization, a composition law, and calibration.
The hierarchy theorem fixes $\varphi$ from self-similar scaling and additive
closure. Two theorems, two hypothesis packages, two constants. Nothing in that
account forces the two to share an origin.

This note shows they do, in a sharp and checkable sense. The shared origin is the
reciprocal involution
\[
  \iota \colon \mathbb{R}_{>0}\to\mathbb{R}_{>0},\qquad \iota(x)=x^{-1}.
\]
This is the duality that exchanges a quantity with its reciprocal: the
dual-recognition symmetry. It is the only structural ingredient we will use. From
it, both the cost and the scale fall out as fixed-point data, and the two
constants turn out to be the two fixed points of a single one-parameter family of
maps built from $\iota$.

The argument is short because the content is structural rather than
computational. What makes it worth stating is not difficulty but identification:
it removes the appearance that $J$ and $\varphi$ are independent posits.

\section{The involution}

\begin{definition}[Reciprocal involution]
$\iota(x)=x^{-1}$ on the positive reals.
\end{definition}

\begin{proposition}\label{prop:inv}
$\iota$ is an involution on $\mathbb{R}_{>0}$: $\iota(\iota(x))=x$ for $x>0$. Its
unique fixed point is the unit: for $x>0$, $\iota(x)=x$ if and only if $x=1$.
\end{proposition}

\begin{proof}
$\iota(\iota(x))=(x^{-1})^{-1}=x$. If $x^{-1}=x$ then $x^2=1$, and the only
positive root is $x=1$; conversely $1^{-1}=1$.
\end{proof}

The fixed point of $\iota$ is the unit. Hold onto this. It is the hinge of the
whole note: the unit is not chosen separately from $\iota$, it is read off $\iota$
as the one place the reciprocal symmetry leaves invariant.

\section{Cost side: $J$ is the $\iota$-symmetric cost, and its zero is $\iota$'s fixed point}

The recognition cost is built directly on $\iota$. Writing it as
\[
  \Jcost(x)=\tfrac12\bigl(x+\iota(x)\bigr)-1
\]
makes the dependence explicit: $J$ is the symmetric average of a quantity and its
reciprocal, recentered so the unit costs nothing.

\begin{proposition}[Cost symmetry]\label{prop:Jsym}
For $x>0$, $\Jcost(\iota(x))=\Jcost(x)$.
\end{proposition}

\begin{proof}
$\Jcost(x^{-1})=\tfrac12(x^{-1}+x)-1=\Jcost(x)$. Equivalently, written as a square,
$\Jcost(x)=(x-1)^2/(2x)$ is manifestly invariant under $x\mapsto x^{-1}$ after
clearing denominators.
\end{proof}

So $J$ is invariant under the reciprocal involution: it is constant on
$\iota$-orbits $\{x,x^{-1}\}$. This is the first hypothesis of the cost-uniqueness
theorem, and it is the structural reason $J$ has the shape it has. The sharper
fact is that $\iota$'s fixed point and $J$'s zero are the same point.

\begin{theorem}[The involution's fixed point is the cost's zero]\label{thm:fixed-zero}
For $x>0$,
\[
  \iota(x)=x \quad\Longleftrightarrow\quad \Jcost(x)=0,
\]
both holding exactly at $x=1$.
\end{theorem}

\begin{proof}
By Proposition~\ref{prop:inv}, $\iota(x)=x$ iff $x=1$. By the squared form
$\Jcost(x)=(x-1)^2/(2x)$, for $x>0$ the cost vanishes iff $x=1$. The two
conditions therefore coincide.
\end{proof}

This is more than ``$J$ is $\iota$-symmetric.'' The symmetry axis of $\iota$, the
locus it fixes, is identical to the null set of $J$, the locus of zero
recognition cost. The unit is simultaneously the only self-reciprocal positive
number and the only costless one. The cost does not introduce a new distinguished
point; it inherits $\iota$'s.

\section{Scale side: $\varphi$ is the unique fixed point of $1+\iota$}

The golden ratio enters through the same involution, shifted. Consider the map
\[
  g(x)=1+\iota(x)=1+x^{-1}.
\]
The shift is by $1$, which is exactly $\iota$'s own fixed point
(Proposition~\ref{prop:inv}). So $g$ is the involution perturbed by the unit it
already singles out. Its fixed-point equation is the self-similarity constraint.

\begin{theorem}[$\varphi$ is the unique fixed point of $1+\iota$ above one]\label{thm:phi}
For $x>1$,
\[
  g(x)=x \quad\Longleftrightarrow\quad x=\varphi,
\]
where $\varphi=(1+\sqrt5)/2$ is the golden ratio.
\end{theorem}

\begin{proof}
For $x>0$, $g(x)=x$ reads $1+x^{-1}=x$. Multiplying by $x$ gives $x+1=x^2$, that
is $x^2=x+1$. The positive roots of $x^2=x+1$ are $\varphi$ and $(1-\sqrt5)/2<0$,
so among $x>1$ the unique solution is $\varphi$. Conversely $\varphi$ satisfies
$\varphi^2=\varphi+1$, hence $\varphi=1+\varphi^{-1}=g(\varphi)$.
\end{proof}

The self-similarity that fixes the scale is not a second principle layered on top
of the cost. It is the fixed-point condition for the same involution, displaced by
the unit. Where $\iota$ fixes the unit, $1+\iota$ fixes $\varphi$.

\section{One involution, two constants}

The two constants now sit in a single picture.

\begin{theorem}[The shared generator]\label{thm:main}
The reciprocal involution $\iota$ determines both fundamental quantities through
fixed points:
\begin{enumerate}[label=\textup{(\roman*)}]
  \item the fixed point of $\iota$ is the unit $1$, which is exactly the unique
    zero of the recognition cost $J$;
  \item the fixed point of the shift $1+\iota$ above one is the golden ratio
    $\varphi$, the scale.
\end{enumerate}
\end{theorem}

\begin{proof}
Immediate from Theorems~\ref{thm:fixed-zero} and~\ref{thm:phi}.
\end{proof}

The unit and the scale, the two numbers the framework is built from, are the two
fixed points of one involution family $\{\iota,\,1+\iota\}$. The family is not
arbitrary either: its two members are $\iota$ and $\iota$ shifted by $\iota$'s own
fixed point, so the whole structure is generated by $\iota$ alone with no second
input. The cost is the $\iota$-symmetric average; $\varphi$ is the fixed point of
the minimal shift of $\iota$ by the unit it fixes.

\paragraph{Machine verification.}
Every statement above is proved in Lean~4 against Mathlib in the module
\texttt{IndisputableMonolith.Foundation.UniversalForcing.ReciprocalGenerator}.
Proposition~\ref{prop:inv} is \texttt{recip\_involutive} and
\texttt{recip\_fixed\_iff}; Proposition~\ref{prop:Jsym} is
\texttt{jcost\_recip\_symmetric}; Theorem~\ref{thm:fixed-zero} is
\texttt{recip\_fixed\_iff\_cost\_zero}; Theorem~\ref{thm:phi} is
\texttt{recipShift\_fixed\_iff} with witness \texttt{phi\_is\_recipShift\_fixed};
Theorem~\ref{thm:main} is \texttt{recip\_pins\_unit\_and\_scale}. The bundle is
recorded as \texttt{ReciprocalGeneratorCert} / \texttt{reciprocalGeneratorCert\_holds}.
A \verb|#print axioms| on these declarations reports no \texttt{sorry} and no
Recognition-Science-specific axiom; the cost-symmetry and golden-ratio facts the
proofs invoke (\texttt{Cost.Jcost\_symm}, \texttt{Cost.Jcost\_eq\_zero\_iff},
\texttt{PhiSupport.phi\_unique\_pos\_root}) are themselves elementary real algebra.

\section{What this does, and does not, claim}

The claim is an identification of origin, and it should be read at exactly that
strength, no more.

It does claim that $J$ and $\varphi$ are not independent posits. Both are fixed by
the single involution $\iota$, the cost by symmetry and the scale by the shifted
fixed point, with the unit serving as the literal hinge between the two: it is
$\iota$'s fixed point, it is $J$'s zero, and it is the amount by which $\iota$ is
shifted to produce the map that fixes $\varphi$.

It does not claim that the entire cost-uniqueness theorem is subsumed. Pinning $J$
to its exact functional form needs more than $\iota$-symmetry: it needs
normalization, the composition law, calibration, and continuity. $\iota$-symmetry
is one hypothesis among those, the one that fixes the symmetry axis; the others
fix the profile along it. This note isolates the shared hypothesis, not the full
derivation.

It does not claim that the full $\varphi$-forcing chain is subsumed either. Deriving
$\varphi$ as the scale of a physical hierarchy needs the additive-closure and
self-similarity content of that hierarchy; what this note shows is that the
\emph{equation} those principles produce, $x=1+x^{-1}$, is the fixed-point equation
of $1+\iota$, so the constant they land on is $\iota$-generated. The physics that
selects this equation is upstream and is not relitigated here.

And it does not claim that physics is reduced to a reciprocal symmetry. The honest
reading is narrower and, for that reason, defensible: two quantities the framework
had treated as separate fundamentals are two fixed points of one involution. That
is a structural unification of inputs, not a theory of everything. It is the kind
of result that makes a framework smaller, which is the direction a foundational
program should want to move.

\end{document}
